\section{Actual direct-image lines and stabilized orders}
\label{inj:stabilized}

This proof of natural invariant transfer retains the actual adjoint
direct-image sheaves throughout. A birational toroidal equidimensional
model makes their invariant rank-one summands line bundles. Local
frames of these lines lift to holomorphic adjoint sections, which will
give the integrability needed for their quotient metrics.

The boundary orders are minima of finitely many affine functions.
After the slopes are made integral, the lines vary affinely in all
large degrees. The intercept comparison also embeds the zero-degree
line, twisted by the eventual slope, into those large-degree lines.
Strong openness then carries strict positivity from degree zero along
this family, preserving the nonzero constant term in Euler interpolation.

\begin{proposition}[Invariant Iitaka conversion by stabilized orders]
\label{inj:stabilized-transfer}
Let $X$ be smooth connected projective and let an algebraic torus
$T$ act on $X$.  Give $L=-K_X$ its natural linearization and suppose
it admits a smooth semipositive metric.  Let $M$ be a linearized
line bundle with $-M$ pseudoeffective.  If there are nonzero
sections $\sigma_j\in H^0(X,M+jL)^T$ for an unbounded set
$J\subset\mathbb Z_{>0}$, then $H^0(X,mL)^T\ne0$ for some
integer $m>0$.
\end{proposition}

\begin{proof}
\smallskip\noindent\textit{Generic adjoint rigidity.}
Write $D(j)=M+jL$ and fix $j_0<j_1$ in $J$.
Choose $k>0$ so that the invariant system
$H^0(X,kD(j_1))^T$ has maximal image dimension among these systems
for all positive degrees.  No finite-generation assertion is needed
to attain this maximum, an integer between zero and $\dim X$.
An equivariant principalization of its base ideal gives a smooth
projective $\pi_0:W_0\to X$, isomorphic off the base locus, and
a connected-fiber Stein map $f_0:W_0\to Y_0$ such that
\[
 \pi_0^*kD(j_1)=f_0^*H+G,
\]
where $H$ is ample on the normal projective $Y_0$ and $G$ is the
effective fixed divisor.  The base action and the fiber action on
$H$ are trivial, while the defining section of $G$ is invariant.
These identifications follow by dividing the invariant generating
sections by the common fixed factor.  The action on the finite
Stein factor is trivial by connectedness of $T$.

Let $K=\mathbb C(Y_0)$ and let $F$ be the generic fiber.  It is
normal and integral, with $H^0(F,\mathcal O_F)=K$.
The field $K$ is relatively algebraically closed in $K(F)$:
an algebraic element has no poles on proper normal $F$, as follows
from a monic equation over $K$.  For every $h>0$,
\begin{equation}\label{inj:invariant-fiber-rigidity}
 \dim_K H^0(F,h\pi_0^*D(j_1)|_F)^T=1.
\end{equation}
There is the section $\pi_0^*\sigma_{j_1}^h$.  If there were
two independent invariants, their ratio would be transcendental
over $K$.  They extend as invariant rational sections on $W_0$
with only vertical poles.  A common pullback of a section of a
sufficiently large multiple $lH$ clears these poles: prescribe
vanishing on their finitely many proper base images.  Multiplying
also by $s_G^l$ produces invariant sections of a multiple of
$D(j_1)$ on $X$.  Products with the original system retain its
image field and add their new ratio, contradicting maximal image
dimension.  This proves \eqref{inj:invariant-fiber-rigidity}.

Let $D_i^F$ be the effective divisor of $\sigma_i|_F$.
The invariant products of the same linearized degree are proportional,
so for $c>j_1$ in $J$,
\[
 (c-j_0)D_{j_1}^F
 =(c-j_1)D_{j_0}^F+(j_1-j_0)D_c^F.
\]
Thus $D_{j_1}^F-D_{j_0}^F\ge0$.  Put $d=j_1-j_0>0$ and
$z=\sigma_{j_1}/\sigma_{j_0}$, a rational invariant section
of $dL$.  It is regular on $F$, with zero support in $D_{j_1}^F$.
If the base is a point we already have a regular invariant section
on $W_0$, hence on $X$.

Otherwise let $e_0$ be the Jacobian section of
$K_{W_0}+\pi_0^*L$.  Its zeros lie over the principalized base
locus and, on $F$, in the support of $D_{j_1}^F$.  For every fixed $s\ge0$, choose $h$ large enough that, on $F$,
\[
 hD_{j_1}^F\ge \operatorname{div}(e_0|_F)
                         +s\operatorname{div}(z|_F).
\]
Such an $h$ exists because both effective divisors on the right
have support in $D_{j_1}^F$.  The invariant quotient
$(\sigma_{j_1}|_F)^h/(e_0z^s|_F)$ is regular.  Multiplication
by this nonzero quotient injects the entire invariant adjoint
section space into the one-dimensional space in
\eqref{inj:invariant-fiber-rigidity}.  Consequently
\begin{equation}\label{inj:adjoint-generic-line}
 H^0\bigl(F,(K_{W_0}+(1+sd)\pi_0^*L)|_F\bigr)^T
       =K\cdot(e_0z^s)|_F\qquad(s\ge0).
\end{equation}

\smallskip\noindent\textit{Actual lines and affine boundary orders.}
The generic rank computation will now be upgraded to line bundles
whose local frames lift to adjoint sections. After further birational
modifications, the same adjoint ranks are preserved by canonical pushforward on a common smooth resolution.
No support claim about new exceptional primes is needed.
Now apply birational weak toroidalization
\cite[Theorem 1.1]{ADK2013}, followed by equidimensionalization
by toroidal subdivisions \cite[Proposition 4.4]{AK2000}.
This gives a birational modification $Y\to Y_0$ with $Y$ smooth,
a normal toroidal $Z$ birational over $W_0$, and an equidimensional
$g:Z\to Y$.  These are birational modifications; no alteration
or assertion of reduced fibers is being used.  Take a smooth
resolution $W\to Z$ and write $f:W\to Y$, $\pi:W\to X$,
and $\pi_Z:Z\to X$.

Toroidal singularities are rational.  Therefore, for
\[
 \mathscr E_s=f_*\mathcal O_W(K_{W/Y}+(1+sd)\pi^*L),
\]
the sheaf $\mathscr E_s\otimes\mathcal O_Y(K_Y)$ equals
\[
 g_*\bigl(\omega_Z\otimes
                  \pi_Z^*\mathcal O_X((1+sd)L)\bigr).
\]
It is reflexive: deleting a codimension-two subset of $Y$ deletes
only codimension at least two from $Z$, by equidimensionality,
and the reflexive canonical sheaf on $Z$ extends its sections.
The pushforward is torsion-free as well, since the source is
integral and the map is dominant.

The toroidal modifications may fail to be equivariant.
Nevertheless $\mathscr E_s$ has a $T$-linearization.  Compute the
same direct image on an equivariant smooth resolution $W'$ of
the main component of $W_0\times_{Y_0}Y$.  A common smooth
resolution of $W$ and $W'$ identifies their canonical direct
images, since all coefficient twists are pulled back from $X$
and a birational morphism of smooth varieties pushes its canonical
sheaf to the canonical sheaf.  The identifications are the
identifications of rational forms, so preserve the generic torus
action.  The invariant part
\[
 A_s=(\mathscr E_s)^T
\]
is now an actual reflexive direct summand of rank one by
\eqref{inj:adjoint-generic-line}; on smooth $Y$ it is a line
bundle.  Birational canonical pushforward gives
\begin{equation}\label{inj:global-invariant-identification}
 H^0(Y,K_Y+A_s)=H^0(X,sdL)^T,
 \qquad h^0(Y,K_Y+A_0)=1.
\end{equation}
The last equality uses the natural invariant trivialization of
$K_X+L$.

We compute these actual lines.  On the regular locus of $Z$ let
$e_Z$ be the Jacobian section with effective divisor $E_Z$.
Use the rational generator $e_Zz^s$ for $K_Y+A_s$.
For each prime $B\subset Y$, let $Q_i$ run over primes of $Z$
mapping onto $B$, and let
\[
 b_i=\operatorname{ord}_{Q_i}(g^*B)>0,
 \quad e_i=\operatorname{ord}_{Q_i}(E_Z),
 \quad c_i=\operatorname{ord}_{Q_i}(z).
\]
A rational base multiplier $a$ is regular above the generic point
of $B$ precisely when
$b_i\operatorname{ord}_B(a)+e_i+sc_i\ge0$ for all $i$.
There are no horizontal poles by
\eqref{inj:adjoint-generic-line}.  The exact allowed pole order is
\[
 u_s(B)=\left\lfloor\min_i\frac{e_i+sc_i}{b_i}\right\rfloor,
 \qquad K_Y+A_s=\mathcal O_Y(C_s),
 \quad C_s=\sum_Bu_s(B)B.
\]
Only finitely many primes occur with a nonzero slope or intercept.
Replace $d,z$ by a common positive multiple and power so that all
$v_B=\min_i(c_i/b_i)$ are integers, and write
$V=\sum_Bv_BB$.  For all sufficiently large integer $s$,
\begin{equation}\label{inj:affine-stabilization}
 C_s=C_\infty+sV,
 \qquad C_\infty\ge C_0.
\end{equation}
Indeed only the indices of minimal slope can minimize for large
$s$.  Among them the least intercept, after taking its floor,
is at least the unrestricted least intercept at zero.  This proves
the stated inequality coefficient by coefficient.  Set
$A_\infty=\mathcal O_Y(C_\infty)-K_Y$.  Then
$A_s=A_\infty+sV$ eventually, and
\begin{equation}\label{inj:line-inclusion}
 A_0+sV\hookrightarrow A_s
\end{equation}
for those $s$. The eventual family has intercept $A_\infty$, but this
inclusion compares it with the actual zero-degree line $A_0$. It is
this inclusion that will transfer sections produced by interpolation
from exponent zero.

\smallskip\noindent\textit{Integrable metrics and interpolation from zero.}
The original moving part, together with $-M$ pseudoeffective,
shows that $\pi^*L-\epsilon f^*H_1$ is pseudoeffective for some
ample $H_1$ on $Y$ and $\epsilon>0$.  The pullback of $H$ to
$Y$ is big, so it still dominates a small ample class.
Mix a metric representing this pseudoeffectivity with the smooth
semipositive metric on $\pi^*L$, using a small positive coefficient.
This gives at $s=0$ a singular metric with trivial multiplier ideal
and curvature at least $\delta f^*\omega_Y$, with $\delta>0$.
For $s>0$ use the smooth semipositive metrics on
$(1+sd)\pi^*L$.

The singular adjoint metric and its extension
\cite[Theorems 3.3.4--3.3.5]{PT2018} give semipositive metrics
on $\mathscr E_s$ on its locally free locus.  The generic
multiplier-ideal inclusion is an isomorphism since the upstairs
ideal is trivial.  Give $A_s$ the quotient metric for the
holomorphic projection onto this summand; quotient positivity is
\cite[Lemma 2.3.4(2)]{PT2018}.  Write its weights as $\psi_s$.
Their description is the logarithm of the squared norm of the
dual line frame in $\mathscr E_s^\vee$.  They extend
plurisubharmonically across the remaining codimension-two subset
of smooth $Y$, since $A_s$ is a line bundle.  They are not
identically $-\infty$, by the nondegenerate finite fiber metric
at almost every smooth base-change parameter.  An invariant
choice of upstairs metric is unnecessary for this quotient argument.

Crucially,
\begin{equation}\label{inj:base-integrability}
 e^{-\psi_s}\in L^1_{\rm loc}(Y)\qquad(s\ge0).
\end{equation}
A local frame of the actual summand lifts to a local section of
$\mathscr E_s$.  Its fiberwise squared norm bounds the quotient
squared norm from above.  The lift is a holomorphic adjoint section
upstairs, after adding a base canonical frame.  Its absolute
density is integrable over inverse images of relatively compact
charts by the trivial upstairs ideal and properness.  Fubini on
the smooth locus identifies this integral with the fiber norms
integrated against the base canonical volume, proving
\eqref{inj:base-integrability}.  A frame of a reflexive hull
without such a lift would not suffice.

The weight $\psi_0$ has a strictly positive lower curvature bound
on all of $Y$.  To check this across the boundary, choose a smooth
function $p$ on $Y$ whose Hessian dominates $\omega_Y$ near a
given point and is bounded above by $C\omega_Y$ globally.
Multiply the upstairs squared metric by $e^{t f^*p}$ for small
$t>0$.  It remains semipositive and has the same ideal.  Its
direct-image quotient weight is $\psi_0-tp$ almost everywhere,
so it is plurisubharmonic everywhere by the extension theorem.
A finite covering gives $i\partial\bar\partial\psi_0
\ge\delta'\omega_Y$ for some $\delta'>0$.

We now put strictly positive metrics with trivial multiplier ideal
on $A_0+sV$ for every $s\ge0$. These are the lines on the left of
\eqref{inj:line-inclusion}, so vanishing for them will retain the
known section count at zero. Fix an integer $s>0$. Take a sufficiently
large integer $u$ in
the range of \eqref{inj:affine-stabilization}, and choose a smooth
metric $\psi_\infty$ on $A_\infty$.  In compatible frames the
formula
\begin{equation}\label{inj:openness-mixture}
 \rho_s=\psi_0+\frac{s}{u}(\psi_u-\psi_\infty)
\end{equation}
defines a metric on $A_0+sV$.  Its curvature is strictly positive
once $s/u$ is small enough.  Strong openness
\cite[Theorem 1.1]{GuanZhou2015} and compactness give
$\eta>0$ such that $e^{-(1+\eta)\psi_0}$ is locally
integrable on a finite cover.  H\"older, together with
\eqref{inj:base-integrability} for this $u$, proves
$e^{-\rho_s}$ integrable if
\[
 \frac{s}{u}\frac{1+\eta}{\eta}\le1.
\]
The smooth $\psi_\infty$ contributes only bounded factors.
No exponent improvement uniform in $u$ is needed.  At $s=0$
the original $\psi_0$ already has these properties.

Nadel vanishing consequently gives
\[
 H^i(Y,K_Y+A_0+sV)=0\qquad(i>0,\ s\ge0).
\]
The Euler polynomial in $s$ has value one at zero by
\eqref{inj:global-invariant-identification}.  Thus it is nonzero
for all sufficiently large integers $s$.  Using
\eqref{inj:line-inclusion} and then
\eqref{inj:global-invariant-identification} gives a nonzero
invariant section of $sdL$ on $X$.
\end{proof}

Under the hypotheses $H^q(X,\mathcal O_X)=0$ for $q>0$,
Proposition~\ref{har:prop:index} and Lemma~\ref{har:lem:forms} supply the
unbounded invariant twisted forms required to apply this proposition.
For positive differential degree, their saturated generic determinant
has pseudoeffective inverse by the cotangent-tensor theorem, exactly
as in Corollary~\ref{inj:forms-input}, now preserving the action.
For degree zero the invariant section is already available.  Thus
this construction is a separate proof of invariant nonvanishing
under those hypotheses, with strong openness and the actual
direct-image line replacing the bounded base metric used earlier.

Section~\ref{sec:descent} applies the invariant conclusions to compact
factors and descends their products through the retained monodromy
and a finite-cover norm.
